%===============================================================================
% DRAFT: Appendix Section - Comparative Analysis of CRL Methods for CRISPRa Data
%===============================================================================
% This draft outlines an appendix section comparing CRL methods and their
% applicability to soft intervention (CRISPRa) settings.
%===============================================================================

\section{Comparative Analysis of CRL Methods for Soft Interventions}
\label{app:crl_methods}
%-------------------------------------------------------------------------------

\subsection{Hard vs. Soft Interventions}
\label{app:crl_methods:intervention_types}
%-------------------------------------------------------------------------------

In the \gls{scm} framework, interventions are characterized by their degree of
modification to the causal structure:

\begin{itemize}
    \item \textbf{Hard interventions} replace the structural equation of a target
          variable entirely, severing all incoming causal edges. Formally, for a
          target variable $\latentcausalvar_i$, a hard intervention sets
          $\latentcausalvar_i = c$ (constant) or $\latentcausalvar_i = \noise_i$,
          eliminating any dependence on its causal parents.
    \item \textbf{Soft interventions} modify only the functional relationship
          between a variable and its parents without breaking causal connections.
          In CRISPRa, activation of a target gene alters its expression level
          while maintaining its regulatory relationships with upstream genes.
\end{itemize}

The distinction is critical: hard interventions provide unambiguous statistical
signatures for causal discovery, while soft interventions preserve some
statistical dependencies, making identifiability fundamentally more challenging.

%-------------------------------------------------------------------------------

\subsection{Theoretical Limits of CRL under Soft Interventions}
\label{app:crl_methods:theory}
%-------------------------------------------------------------------------------

\subsubsection{Linear Models: Ancestor Mixing}
\citet{squires2023} established that under soft interventions, linear latent models
are identifiable only up to ancestral relations. Formally, let $\latentcausalvars = \因果$
be generated by a linear \gls{scm} with soft intervention on target $i$:

\begin{equation}
    \latentcausalvar_i^{I} = \gamma_i + \sum_{j \in \parents(i)} \beta_{ij} \latentcausalvar_j + \noise_i'
\end{equation}

where $\gamma_i$ is the intervention shift. The identifiability result shows that
the true causal structure $\graph$ can only be recovered up to ancestors of
intervened variables, not the full DAG.

\subsubsection{Multi-Node Interventions}
\citet{varici2024multi} extend this analysis to multi-node interventions,
showing that perfect identifiability requires hard interventions. With soft
interventions on multiple nodes, the identifiability degrades to:

\begin{itemize}
    \item For single soft intervention: ID up to ancestors of target
    \item For multi-node soft interventions: Entanglement with direct parents
\end{itemize}

\subsubsection{Dense Graphs and Polytrees}
\citet{zhang2023identifiability} prove perfect identifiability under soft
interventions but require a \emph{polytree} structure (sparse, directed acyclic
graph without cycles). Biological gene regulatory networks are inherently dense,
violating this assumption and leading to ancestor mixing.

%-------------------------------------------------------------------------------

\subsection{Method-by-Method Analysis}
\label{app:crl_methods:methods}
%-------------------------------------------------------------------------------

\subsubsection{Squires et al. (ICML 2023)}
\textit{Transform:} Linear. \textit{Identifiability:} Partial (ID up to ancestors).
\citet{squires2023} provide a foundational result showing that without hard
interventions, perfect causal discovery is impossible. The model cannot
distinguish a CRISPRa-perturbed gene from its upstream ancestors.

\subsubsection{Buchholz et al. (NeurIPS 2023)}
\textit{Transform:} Linear Gaussian + Nonparametric. \textit{Identifiability:}
Perfect for soft shifts under nonlinear mixing. However, theory depends on
sparsity assumptions that fail in dense biological networks.

\subsubsection{Varici et al. (JMLR 2025, NeurIPS 2024)}
\textit{Transform:} Linear / General. \textit{Identifiability:} Partial.
These works explicitly guarantee only partial identifiability under soft
interventions. Direct biological parents remain entangled with targets.

\subsubsection{Zhang et al. (NeurIPS 2023)}
\textit{Transform:} Polynomial (Nonlinear). \textit{Identifiability:} Perfect
for polytree structures. Best theoretical fit among existing methods, but
biological GRNs are dense, causing ancestor mixing.

\subsubsection{Brehmer et al. (NeurIPS 2022)}
\textit{Transform:} Nonparametric. \textit{Identifiability:} Conditional on
paired observations. Requires weakly supervised paired data (same cell before/after
perturbation), which is physically impossible in scRNA-seq due to cellular lysing.

\subsubsection{SoftILCM (arXiv 2024)}
\textit{Transform:} Implicit (VAE). \textit{Identifiability:} Partial.
Disentanglement bounds degrade significantly with increasing DAG density and
weakening intervention strength---both characteristics of CRISPRa experiments.

%-------------------------------------------------------------------------------

\subsection{Implications for \gls{raptorgraph}}
\label{app:crl_methods:raptorgraph}
%-------------------------------------------------------------------------------

\gls{raptorgraph} addresses these theoretical gaps through:

\begin{enumerate}
    \item \textbf{Preconditioned GraphPathway Encoder:} Enforces sparse,
          one-to-one mappings from perturbed genes to causal latent factors,
          mitigating the soft intervention ambiguity.
    \item \textbf{Optimal Transport Alignment:} Accounts for distribution shift
          without requiring paired observations, bypassing the data modality
          limitation.
    \item \textbf{DAGMA-Based Causal Structure Learning:} Learns the latent
          causal graph directly from data, accommodating dense biological
          networks without requiring polytree assumptions.
\end{enumerate}

%===============================================================================
% END OF DRAFT
%===============================================================================
